% Detalhar os conceitos a ser abordado no artigo%

Nesta seção, são apresentados trabalhos relacionados ao uso de Modelos de Linguagem em contextos agrícolas e na pecuária leiteira, destacando suas contribuições, limitações e a forma como dialogam com os objetivos desta pesquisa.

 O \textit{DairyGPT} foi um estudo de \cite{GONTIJO2025101097}, o qual investiga a interação entre produtores de leite e modelos de LLMs treinados exclusivamente com dados numéricos, explorando como essa abordagem pode contribuir para a transformação digital do setor. Os autores demonstram que a especialização do modelo com dados da pecuária leiteira aumenta a precisão e a relevância das respostas geradas. Para isso, aplicam técnicas como RAG e NL2SQL (\textit{Natural Language to SQL}), permitindo consultas mais eficientes a bancos de dados estruturados. Na etapa de avaliação, o modelo apresentou resultados promissores, indicando que o uso de dados específicos potencializa a aplicabilidade de LLMs em cenários agropecuários. Em contraste, este trabalho realiza uma avaliação da qualidade semântica e da similaridade do conteúdo textual gerado. Enquanto o \textit{DairyGPT} concentra-se no desempenho de consultas a bancos de dados puramente numéricos, fornecendo respostas baseadas na análise de curvas de lactação.

O trabalho de \cite{awais2025agrogptefficientagriculturalvisionlanguage}, denominado AgroGPT, propõe o desenvolvimento de um modelo especializado para o domínio agrícola, com o objetivo de superar as limitações de Modelos Multimodais de Linguagem (LMMs) generalistas, como GPT-4 \cite{openai2023chatgpt}, BARD \cite{Bard} e LLaVA \cite{LLaVA}. Os autores argumentam que, mesmo com técnicas de engenharia de prompt, esses modelos apresentam lacunas de conhecimento no setor agrícola, comprometendo a qualidade das respostas geradas. Como solução, sugerem ajuste fino com dados específicos do domínio, combinado a estratégias de construção de prompts mais detalhados, para aprimorar a eficácia dos modelos. Este estudo reforça a relevância de adaptar LLMs a contextos específicos, evidenciando que a especialização é determinante para o sucesso desses modelos. Entretanto, o foco do AgroGPT concentra-se na análise de imagens de plantas e folhas, utilizando \textit{Expert Tuning} (Tunagem Experiente) e avaliação manual por especialistas, enquanto o presente trabalho direciona-se à análise da qualidade semântica de textos gerados no contexto da pecuária leiteira.

Um trabalho precursor a este é o de \cite{agents_in_dairy_science}, cuja proposta consiste na utilização de modelos LLM da família LLAMA \cite{touvron2023llamaopenefficientfoundation} aprimorados com RAG. A abordagem apresentada pelos autores é dividida em duas partes: a primeira consiste no desenvolvimento de um chatbot auxiliar para apoio à tomada de decisões, fundamentado em artigos provenientes do \textit{Journal of Dairy Science} (JDS), utilizados como base de conhecimento no sistema RAG; a segunda envolve a integração de uma interface em linguagem natural com modelos acadêmicos capazes de visualizar e interpretar predições de resultados. Nesse contexto, o presente trabalho busca complementar esse desenvolvimento inicial ao avaliar outras estratégias de aprimoramento de modelos, como \textit{Zero-shot} e diferentes configurações de sistemas RAG. Além disso, propõe-se analisar a qualidade das respostas geradas por meio de métricas de avaliação textual.

De forma semelhante, o artigo de \cite{zhang-etal-2025-beefbot} apresenta o \textit{Beef Bot}, um \textit{chatbot} voltado para a produção de carne bovina, aprimorado com RAG e \textit{Fine Turning} (Tunagem Fina) \cite{fine_turning_llm_for_specialized_use_cases} para gerar respostas contextualizadas e úteis aos produtores. Enquanto o \textit{Beef Bot} concentra-se na pecuária de corte, este trabalho foca no setor leiteiro, comparando como diferentes modelos podem ser avaliados e aprimorados usando as mesmas técnicas.

No campo da saúde, o estudo de \cite{ChatGPTInMedicice} analisa vantagens e limitações do uso de LLMs, como o \textit{ChatGPT}, na medicina. Entre os benefícios, destacam-se buscas rápidas de artigos, apoio à educação médica e monitoramento de dados de pacientes. Entre os problemas, incluem-se riscos de violação de direitos autorais, manuseio de informações sigilosas e fenômeno de “alucinação”, no qual o modelo gera informações irrelevantes. Este estudo contrasta com a proposta do presente trabalho, pois aqui defendemos a aplicação de técnicas como RAG para aproveitar os dados especializados, garantindo confidencialidade e precisão.

A avaliação da qualidade textual também se mostra essencial em pesquisas com LLMs. O estudo de \cite{nguyen2024comparative} apresenta métricas automáticas amplamente utilizadas para comparação com textos humanos, ressaltando que a avaliação manual, embora mais precisa, é demorada e pouco escalável. Esse ponto reforça a importância de métodos sistemáticos de avaliação para garantir coerência e relevância semântica nos textos gerados.

Neste trabalho, foram selecionadas BLEU e ROUGE-1 para medir similaridade léxica, e BERTScore, BARTScore e GPTScore para avaliar coerência e similaridade semântica, capturando nuances contextuais que métricas tradicionais podem não detectar. Essa combinação permite uma avaliação objetiva e abrangente, reduzindo a necessidade de análises manuais e assegurando que os textos gerados sejam consistentes em relação às referências.

Em síntese, os trabalhos citados evidenciam o desenvolvimento e aprimoramento de modelos voltados para domínios específicos, suas limitações e oportunidades de melhoria. Neste contexto, este trabalho se diferencia por aplicar métricas avançadas de avaliação textual, permitindo uma análise objetiva da qualidade das respostas geradas. Essa abordagem complementa a adaptação dos modelos, direcionada à melhoria de desempenho em um domínio específico, conforme detalhado na próxima seção.
